\scanpage{73}
\mspage{66}
% archive sheet number: 72
S'il ne l'était, par $(\alpha)$ il existerait $I_1\subset I_0$, $I_1\ne I_0$ tel que $I_1$ non admissible. De $=$, il y en aurait $I_2\subset I_1$, $I_2\ne I_1$, non admissible etc. --- on pourrait construire une suite \inserted{str.} décroissante d'ouverts $I_n$ de $I$, contrairement au caractère noethérien de $I^\circ$.

Prouvons $(\alpha)$, et distinguons encore deux cas. a) $I^\circ$ est irréductible, donc admet un pt générique i.e. un plus petit élément, i.e. $I$ admet un plus grand élément. Ce cas a déjà été traité \unsure{tantôt} --- \deleted{toute} \unsure{tte} cat. avec objet final est admissible. b) $I^\circ$ est réductible, donc \deleted{$I^\circ$} réunion de deux fermés $I_1^\circ$, $I_2$ distincts de $I^\circ$, donc $I$ est réunion de deux ouverts $I_1,I_2$ distincts de $I$. Alors par l'hyp. \inserted{de} $(\alpha)$, $I_1,I_2$ \deleted{\unsure{et}} et $I_0=I_1\cap I_2$ sont admissibles. On termine alors comme tantôt, étant ramené à l'admissibilité de
\diagram{
\begin{tikzcd}[ampersand replacement=\&,column sep=small,row sep=small]
\&\&1\\
\unsure{\Psi|\Upsilon}=\&0\arrow[ur]\arrow[dr]\&\\
\&\&2
\end{tikzcd}
}
!
